\batchmode
\documentclass[12pt]{article}
\makeatletter
\usepackage{epsf,amsmath}

\oddsidemargin0in
\textwidth6.5in
\topmargin-0.5in
\textheight9in
\parindent0in

\providecommand{\OneFigure}[1]{\begin{center}
\parbox{1.5in}{\epsfysize=1.5in \epsfbox{#1}}\end{center}}
\par\usepackage[dvips]{color}
\pagecolor[gray]{.7}


\makeatletter
\count@=\the\catcode`\_ \catcode`\_=8 
\newenvironment{tex2html_wrap}{}{} \catcode`\_=\count@
\makeatother
\let\mathon=$
\let\mathoff=$
\ifx\AtBeginDocument\undefined \newcommand{\AtBeginDocument}[1]{}\fi
\newbox\sizebox
\setlength{\hoffset}{0pt}\setlength{\voffset}{0pt}
\addtolength{\textheight}{\footskip}\setlength{\footskip}{0pt}
\addtolength{\textheight}{\topmargin}\setlength{\topmargin}{0pt}
\addtolength{\textheight}{\headheight}\setlength{\headheight}{0pt}
\addtolength{\textheight}{\headsep}\setlength{\headsep}{0pt}
\setlength{\textwidth}{349pt}
\newwrite\lthtmlwrite
\makeatletter
\let\realnormalsize=\normalsize
\global\topskip=2sp
\def\preveqno{}\let\real@float=\@float \let\realend@float=\end@float
\def\@float{\let\@savefreelist\@freelist\real@float}
\def\end@float{\realend@float\global\let\@freelist\@savefreelist}
\let\real@dbflt=\@dbflt \let\end@dblfloat=\end@float
\let\@largefloatcheck=\relax
\def\@dbflt{\let\@savefreelist\@freelist\real@dbflt}
\def\adjustnormalsize{\def\normalsize{\mathsurround=0pt \realnormalsize
 \parindent=0pt\abovedisplayskip=0pt\belowdisplayskip=0pt}\normalsize}%
\def\lthtmltypeout#1{{\let\protect\string\immediate\write\lthtmlwrite{#1}}}%
\newcommand\lthtmlhboxmathA{\adjustnormalsize\setbox\sizebox=\hbox\bgroup}%
\newcommand\lthtmlvboxmathA{\adjustnormalsize\setbox\sizebox=\vbox\bgroup%
 \let\ifinner=\iffalse }%
\newcommand\lthtmlboxmathZ{\@next\next\@currlist{}{\def\next{\voidb@x}}%
 \expandafter\box\next\egroup}%
\newcommand\lthtmlmathtype[1]{\def\lthtmlmathenv{#1}}%
\newcommand\lthtmllogmath{\lthtmltypeout{l2hSize %
:\lthtmlmathenv:\the\ht\sizebox::\the\dp\sizebox::\the\wd\sizebox.\preveqno}}%
\newcommand\lthtmlfigureA[1]{\let\@savefreelist\@freelist
       \lthtmlmathtype{#1}\lthtmlvboxmathA}%
\newcommand\lthtmlfigureZ{\lthtmlboxmathZ\lthtmllogmath\copy\sizebox
       \global\let\@freelist\@savefreelist}%
\newcommand\lthtmldisplayA[1]{\lthtmlmathtype{#1}\lthtmlvboxmathA}%
\newcommand\lthtmldisplayB[1]{\edef\preveqno{(\theequation)}%
  \lthtmldisplayA{#1}\let\@eqnnum\relax}%
\newcommand\lthtmldisplayZ{\lthtmlboxmathZ\lthtmllogmath\lthtmlsetmath}%
\newcommand\lthtmlinlinemathA[1]{\lthtmlmathtype{#1}\lthtmlhboxmathA  \vrule height1.5ex width0pt }%
\newcommand\lthtmlinlineA[1]{\lthtmlmathtype{#1}\lthtmlhboxmathA}%
\newcommand\lthtmlinlineZ{\egroup\expandafter\ifdim\dp\sizebox>0pt %
  \expandafter\centerinlinemath\fi\lthtmllogmath\lthtmlsetinline}
\newcommand\lthtmlinlinemathZ{\egroup\expandafter\ifdim\dp\sizebox>0pt %
  \expandafter\centerinlinemath\fi\lthtmllogmath\lthtmlsetmath}
\def\lthtmlsetinline{\hbox{\vrule width.1em\vtop{\vbox{%
  \kern.1em\copy\sizebox}\ifdim\dp\sizebox>0pt\kern.1em\else\kern.3pt\fi
  \ifdim\hsize>\wd\sizebox \hrule depth1pt\fi}}}
\def\lthtmlsetmath{\hbox{\vrule width.1em\vtop{\vbox{%
  \kern.1em\kern0.8 pt\hbox{\hglue.17em\copy\sizebox\hglue0.8 pt}}\kern.3pt%
  \ifdim\dp\sizebox>0pt\kern.1em\fi \kern0.8 pt%
  \ifdim\hsize>\wd\sizebox \hrule depth1pt\fi}}}
\def\centerinlinemath{%\dimen1=\ht\sizebox
  \dimen1=\ifdim\ht\sizebox<\dp\sizebox \dp\sizebox\else\ht\sizebox\fi
  \advance\dimen1by.5pt \vrule width0pt height\dimen1 depth\dimen1 
 \dp\sizebox=\dimen1\ht\sizebox=\dimen1\relax}

\def\lthtmlcheckvsize{\ifdim\ht\sizebox<\vsize\expandafter\vfill
  \else\expandafter\vss\fi}%
\makeatletter \tracingstats = 1 


\begin{document}
\pagestyle{empty}\thispagestyle{empty}%
\lthtmltypeout{latex2htmlLength hsize=\the\hsize}%
\lthtmltypeout{latex2htmlLength vsize=\the\vsize}%
\lthtmltypeout{latex2htmlLength hoffset=\the\hoffset}%
\lthtmltypeout{latex2htmlLength voffset=\the\voffset}%
\lthtmltypeout{latex2htmlLength topmargin=\the\topmargin}%
\lthtmltypeout{latex2htmlLength topskip=\the\topskip}%
\lthtmltypeout{latex2htmlLength headheight=\the\headheight}%
\lthtmltypeout{latex2htmlLength headsep=\the\headsep}%
\lthtmltypeout{latex2htmlLength parskip=\the\parskip}%
\lthtmltypeout{latex2htmlLength oddsidemargin=\the\oddsidemargin}%
\makeatletter
\if@twoside\lthtmltypeout{latex2htmlLength evensidemargin=\the\evensidemargin}%
\else\lthtmltypeout{latex2htmlLength evensidemargin=\the\oddsidemargin}\fi%
\makeatother
{\newpage\clearpage
\lthtmlinlinemathA{tex2html_wrap_inline216}%
$n \times n$%
\lthtmlinlinemathZ
\hfill\lthtmlcheckvsize\clearpage}

{\newpage\clearpage
\lthtmldisplayA{displaymath164}%
\begin{displaymath}Ax = 0 \ \ \ \ Ay = 0
\end{displaymath}%
\lthtmldisplayZ
\hfill\lthtmlcheckvsize\clearpage}

{\newpage\clearpage
\lthtmldisplayA{displaymath165}%
\begin{displaymath}A(x+y) = Ax + Ay = \v 0 + \v 0 = \v 0.
\end{displaymath}%
\lthtmldisplayZ
\hfill\lthtmlcheckvsize\clearpage}

{\newpage\clearpage
\lthtmldisplayA{displaymath166}%
\begin{displaymath}A(cx) = cAx = c\v 0 = \v 0.
\end{displaymath}%
\lthtmldisplayZ
\hfill\lthtmlcheckvsize\clearpage}

{\newpage\clearpage
\lthtmldisplayA{displaymath167}%
\begin{displaymath}u = 
\left[
\begin{matrix}
1 \\
-1  \\
1
\end{matrix}
\right], \ \ \ \ 
v = 
\left[
\begin{matrix}
2 \\
1  \\
1
\end{matrix}
\right].
\end{displaymath}%
\lthtmldisplayZ
\hfill\lthtmlcheckvsize\clearpage}

{\newpage\clearpage
\lthtmldisplayA{displaymath168}%
\begin{displaymath}u \cdot v = \cos \theta \|u\| \|v\|.
\end{displaymath}%
\lthtmldisplayZ
\hfill\lthtmlcheckvsize\clearpage}

{\newpage\clearpage
\lthtmlinlinemathA{tex2html_wrap_inline238}%
$\arccos(\sqrt{2}/3)$%
\lthtmlinlinemathZ
\hfill\lthtmlcheckvsize\clearpage}

{\newpage\clearpage
\lthtmldisplayA{displaymath169}%
\begin{displaymath}A = \left[
\begin{matrix}
4 & 1 & 2 & 4 \\
-1 & 0 & 1 & 3 \\
10 & 2 & 2 & 2 \\
3 & 1 & 3 & 7
\end{matrix}
\right].
\end{displaymath}%
\lthtmldisplayZ
\hfill\lthtmlcheckvsize\clearpage}

{\newpage\clearpage
\lthtmldisplayA{displaymath170}%
\begin{displaymath}\left[
\begin{matrix}
1 & 0 & -1 & -3 \\
0 & 1 & 6 & 16\\
0 & 0 & 0 & 0 \\
0 & 0 & 0 & 0
\end{matrix}
\right].
\end{displaymath}%
\lthtmldisplayZ
\hfill\lthtmlcheckvsize\clearpage}

{\newpage\clearpage
\lthtmldisplayA{displaymath171}%
\begin{displaymath}\left[
\begin{matrix}
s + 3 t \\
-6 s - 16 t \\
s \\
t
\end{matrix}
\right].
\end{displaymath}%
\lthtmldisplayZ
\hfill\lthtmlcheckvsize\clearpage}

{\newpage\clearpage
\lthtmldisplayA{displaymath172}%
\begin{displaymath}\left\{
\left[
\begin{matrix}
1 \\-6 \\1 \\0
\end{matrix}
\right],
\left[
\begin{matrix}
3 \\-16 \\0 \\1
\end{matrix}
\right]
\right\}.
\end{displaymath}%
\lthtmldisplayZ
\hfill\lthtmlcheckvsize\clearpage}

{\newpage\clearpage
\lthtmldisplayA{displaymath174}%
\begin{displaymath}A =
\left[
\begin{matrix}
3 & 0 & 0 \\
0 & 2 & 0 \\
0 & 0 & 1/2
\end{matrix}
\right].
\end{displaymath}%
\lthtmldisplayZ
\hfill\lthtmlcheckvsize\clearpage}

{\newpage\clearpage
\lthtmldisplayA{displaymath175}%
\begin{displaymath}u^T A v = 3 u_1 v_1 + 2u_2 v_2 + \frac{1}{2} u_3 v_3.
\end{displaymath}%
\lthtmldisplayZ
\hfill\lthtmlcheckvsize\clearpage}

{\newpage\clearpage
\lthtmlinlinemathA{tex2html_wrap_inline260}%
$(u,u) \geq 0$%
\lthtmlinlinemathZ
\hfill\lthtmlcheckvsize\clearpage}

{\newpage\clearpage
\lthtmlinlinemathA{tex2html_wrap_inline264}%
$u = \v 0$%
\lthtmlinlinemathZ
\hfill\lthtmlcheckvsize\clearpage}

{\newpage\clearpage
\lthtmldisplayA{displaymath177}%
\begin{displaymath}\|u+v\|^2 = \|u\|^2 + \|v\|^2
\end{displaymath}%
\lthtmldisplayZ
\hfill\lthtmlcheckvsize\clearpage}

{\newpage\clearpage
\lthtmldisplayA{eqnarraystar46}%
\begin{eqnarray*}\|u+v\|^2 & = & (u+v,u+v) \\
& = & (u,u) + 2(u,v) + (v,v) \\
& = & \|u\|^2 + \|v\|^2 + 2(u,v)
\end{eqnarray*}%
\lthtmldisplayZ
\hfill\lthtmlcheckvsize\clearpage}

{\newpage\clearpage
\lthtmldisplayA{displaymath178}%
\begin{displaymath}W = {\rm span}
\left\{
\left[
\begin{matrix}
3  \\
0  \\
0 \\
1
\end{matrix}
\right],
\left[
\begin{matrix}
1  \\
0 \\
1  \\
1
\end{matrix}
\right],
\left[
\begin{matrix}
0 \\
-1  \\
0  \\
1
\end{matrix}
\right]
\right\}.
\end{displaymath}%
\lthtmldisplayZ
\hfill\lthtmlcheckvsize\clearpage}

{\newpage\clearpage
\lthtmldisplayA{displaymath179}%
\begin{displaymath}\overline{v}_1 = 
\left[
\begin{matrix}
3  \\
0  \\
0 \\
1
\end{matrix}
\right]
\end{displaymath}%
\lthtmldisplayZ
\hfill\lthtmlcheckvsize\clearpage}

{\newpage\clearpage
\lthtmldisplayA{displaymath180}%
\begin{displaymath}\overline{v}_2 = w_2 - 
\frac{w_2 \cdot \overline{v}_1}{\|\overline{v}_1\|^2}\overline{v}_1 =
\left[
\begin{matrix}
-\frac{1}{5}  \\
0  \\
1 \\
\frac{3}{5}
\end{matrix}
\right]
\end{displaymath}%
\lthtmldisplayZ
\hfill\lthtmlcheckvsize\clearpage}

{\newpage\clearpage
\lthtmldisplayA{displaymath181}%
\begin{displaymath}\overline{v}_3 = w_3 - 
\frac{w_3 \cdot \overline{v}_1}{\|\overline{v}_1\|^2}\overline{v}_1 -
\frac{w_3 \cdot \overline{v}_2}{\|\overline{v}_2\|^2}\overline{v}_2 =
\left[
\begin{matrix}
-\frac{3}{14}  \\
-1  \\
-\frac{3}{7} \\
\frac{9}{14}
\end{matrix}
\right]
\end{displaymath}%
\lthtmldisplayZ
\hfill\lthtmlcheckvsize\clearpage}

{\newpage\clearpage
\lthtmldisplayA{displaymath182}%
\begin{displaymath}\left\{
\frac{\sqrt{10}}{10}
\left[
\begin{matrix}
3  \\
0  \\
0 \\
1
\end{matrix}
\right]
,
\sqrt{\frac{5}{7}}
\left[
\begin{matrix}
-\frac{1}{5}  \\
0  \\
1 \\
\frac{3}{5}
\end{matrix}
\right]
,
\sqrt{\frac{14}{23}}
\left[
\begin{matrix}
-\frac{3}{14}  \\
-1  \\
-\frac{3}{7} \\
\frac{9}{14}
\end{matrix}
\right]
\right\}
\end{displaymath}%
\lthtmldisplayZ
\hfill\lthtmlcheckvsize\clearpage}

{\newpage\clearpage
\lthtmlinlinemathA{tex2html_wrap_inline302}%
$L: P^3 \rightarrow P^3$%
\lthtmlinlinemathZ
\hfill\lthtmlcheckvsize\clearpage}

{\newpage\clearpage
\lthtmldisplayA{displaymath183}%
\begin{displaymath}L(p) = \frac{1}{2} x \frac{dp}{dx} - p.
\end{displaymath}%
\lthtmldisplayZ
\hfill\lthtmlcheckvsize\clearpage}

{\newpage\clearpage
\lthtmlinlinemathA{tex2html_wrap_inline304}%
$\ker L$%
\lthtmlinlinemathZ
\hfill\lthtmlcheckvsize\clearpage}

{\newpage\clearpage
\lthtmldisplayA{displaymath185}%
\begin{displaymath}L(p) = \frac{1}{2} a x^3 -\frac{1}{2} c x -d
\end{displaymath}%
\lthtmldisplayZ
\hfill\lthtmlcheckvsize\clearpage}

{\newpage\clearpage
\lthtmldisplayA{displaymath186}%
\begin{displaymath}S = \left\{ 
\left[
\begin{matrix}
1 \\3 \\2
\end{matrix}
\right],
\left[
\begin{matrix}
1 \\-1 \\0
\end{matrix}
\right],
\left[
\begin{matrix}
2 \\1 \\1
\end{matrix}
\right]
\right\}, \ \ \ \ \
T = 
\left\{ 
\left[
\begin{matrix}
1 \\-1 \\0
\end{matrix}
\right],
\left[
\begin{matrix}
1 \\0 \\0
\end{matrix}
\right],
\left[
\begin{matrix}
0 \\-1 \\1
\end{matrix}
\right]
\right\}.
\end{displaymath}%
\lthtmldisplayZ
\hfill\lthtmlcheckvsize\clearpage}

{\newpage\clearpage
\lthtmlinlinemathA{tex2html_wrap_inline328}%
$P_{S \leftarrow T}$%
\lthtmlinlinemathZ
\hfill\lthtmlcheckvsize\clearpage}

{\newpage\clearpage
\lthtmldisplayA{displaymath187}%
\begin{displaymath}[u]_T =
\left[
\begin{matrix}
2 \\0 \\1
\end{matrix}
\right].
\end{displaymath}%
\lthtmldisplayZ
\hfill\lthtmlcheckvsize\clearpage}

{\newpage\clearpage
\lthtmldisplayA{displaymath188}%
\begin{displaymath}P_{S \leftarrow T} = 
\left[
\begin{matrix}
0 & -\frac{1}{2} & 2 \\
1 & -\frac{1}{2} & 4 \\
0 & 1 & -3 \\
\end{matrix}
\right].
\end{displaymath}%
\lthtmldisplayZ
\hfill\lthtmlcheckvsize\clearpage}

{\newpage\clearpage
\lthtmldisplayA{displaymath189}%
\begin{displaymath}[u]_S = \left[
\begin{matrix}
2 \\6 \\-3
\end{matrix}
\right].
\end{displaymath}%
\lthtmldisplayZ
\hfill\lthtmlcheckvsize\clearpage}


\end{document}
